\section[F. 迷路日]{\myboxvi{      \BF{\TT{F}}} 迷路日}

\textit{运行时间限制}   1 second

\textit{内存使用限制} 1 Gigabyte

\textit{读入文件} \qquad 标准输入

\textit{输出文件} \qquad 标准输出

\begin{quote}
   \textit{月が話しかけてきた\\
そして不思議な力をくれた\\}
   \textit{也许该交换一下标题。}
\end{quote}

\subsection{题目描述}

小 T 来到了 T 城，她要前往一个 Livehouse 演出。

T 城可以用一个无限大的欧氏二维平面表示，$y$ 轴正方向表示正北 (N) 方向，$x$ 轴正方向表示正东 (E) 方向，$y$ 轴负方向表示正南 (S) 方向，$x$ 轴负方向表示正西 (W) 方向。

小 T 初始站在 $(x_{t}, y_{t})$ 的位置，初始面向正东、正南、正西、正北四个方向中的一个，她要去的 Livehouse 位于 $(x_l, y_l)$.

她每次移动可以选择下面两种中的一个：

\begin{itemize}
   \item 直行，沿当前方向走一个单位长度。
   \item 右转，将小 T 面对的方向顺时针转 $90^\circ$ （北转向东，东转向南，南转向西，西转向北）。
\end{itemize}

注意小 T 不能左转或掉头。

请求出在上述条件下，小 T 至少要\textbf{右转}几次才能到 Livehouse, 可以证明在这些规则下总能到达目标位置。

\subsection{输入格式}

第一行输入一个整数 $T$ ($1 \le T \le 10^5$), 表示测试数据的数量。

每组测试数据一行输入 5 个元素 $x_t$, $y_t$, $dir$, $x_l$, $y_l$($x_t, y_t, x_l, y_l \in \{n \in \mathbb{Z} : -10^9 \le n \le 10^9$ \}, $dir \in \{\text{N}, \text{E}, \text{S}, \text{W}\}$), 分别表示小 T 所在的位置，小 T 初始面对的方向，Livehouse 的位置。

\subsection{输出格式}

对于每组测试数据，输出一行一个整数，表示在上述条件下，小 T 至少要\textbf{右转}几次才能到 Livehouse.

\subsection{样例 1}

\begin{minipage}{\linewidth}
   \begin{minipage}[t]{0.49\linewidth}
      \begin{lstlisting}[title=\texttt{standard input}]
8
0 0 E 3 -2
1 1 N 0 1
10 10 E 10 10
-9 2 E 3 2
0 -7 E -1 -4
0 6 E -4 6
10 -9 S 10 9
1 -3 N -1 0
\end{lstlisting}
   \end{minipage}
   \begin{minipage}[t]{0.49\linewidth}
      \begin{lstlisting}[title=\texttt{standard output}]
1
3
0
0
3
2
2
3
\end{lstlisting}
   \end{minipage}
\end{minipage}

\subsection{样例 1 解释}

对于第 $1$ 组测试数据，首先向东走 $3$ 步到达 $(3,0)$, 随后右转朝向南，再向南走 $2$ 步到 $(3, -2)$.

对于第 $2$ 组测试数据，首先右转三次朝向西，随后走 $1$ 步到达 $(0, 1)$.

对于第 $3$ 组测试数据，起点和终点重合因此不需要右转。